\documentclass[11pt]{article}

% ----------------------------------------------------------------------
%  The Synergy Reality Test — Practitioner's Guide (methods companion)
%  The "show me how" document behind synergy-reality-test.tex.
%  Draft prepared by the ENT 105 team for Creo Advisors.
%  Compile: pdflatex synergy-reality-test-guide.tex  (run twice for ToC)
% ----------------------------------------------------------------------

\usepackage[utf8]{inputenc}
\usepackage[T1]{fontenc}
\usepackage{textcomp}
\usepackage[margin=1in]{geometry}
\usepackage{booktabs}
\usepackage{array}
\usepackage{amsmath}
\usepackage{amssymb}
\usepackage{enumitem}
\usepackage{titlesec}
\usepackage{xcolor}
\usepackage{tcolorbox}
\usepackage[hidelinks]{hyperref}

\definecolor{creonavy}{RGB}{17,38,73}
\definecolor{creoaccent}{RGB}{27,94,140}

\titleformat{\section}{\large\bfseries\color{creonavy}}{\thesection}{0.6em}{}
\titleformat{\subsection}{\normalsize\bfseries\color{creonavy!85!black}}{\thesubsection}{0.5em}{}
\titlespacing*{\section}{0pt}{1.6ex plus .3ex}{0.7ex}
\setlist[itemize]{leftmargin=1.3em,topsep=2pt,itemsep=2.5pt}
\setlist[enumerate]{leftmargin=1.6em,topsep=2pt,itemsep=3pt}
\setlength{\parskip}{4pt plus 1pt minus 1pt}
\setlength{\parindent}{0pt}
\setlength{\emergencystretch}{2em}

\newtcolorbox{note}{colback=creoaccent!7,colframe=creoaccent,
  boxrule=0.5pt,arc=2pt,left=7pt,right=7pt,top=5pt,bottom=5pt}
\newtcolorbox{judgment}{colback=orange!7,colframe=orange!75!black,
  boxrule=0.5pt,arc=2pt,left=7pt,right=7pt,top=5pt,bottom=5pt}

\newcommand{\durable}{\textcolor{creoaccent}{\textbf{durable}}}
\newcommand{\fragile}{\textcolor{gray!55!black}{\textbf{fragile}}}

\title{\textbf{\color{creonavy} The Synergy Reality Test --- Practitioner's Guide}\\[3pt]
\large How to run the analysis: the rent split, the entry floor, worked end-to-end}
\author{\normalsize Methods companion to the one-pager --- for analysts \textbf{(internal)}}
\date{\normalsize 2 June 2026}

\begin{document}
\maketitle
\vspace{-1.5em}

\begin{note}
\textbf{What this is.} The one-pager \emph{sells} the Synergy Reality Test; this guide
\emph{runs} it. It is what you reach for when a client (or a skeptical colleague) asks ``show
me how you split durable from fragile and where the floor comes from.'' Same roll-up as the
one-pager, carried through end to end.
\end{note}

\tableofcontents
\vspace{0.6em}\hrule\vspace{0.8em}

% ======================================================================
\section{The principle: a roll-up captures rent, it does not create it}
% ======================================================================
The governing fact is conservation. The total economic rent in a market is fixed;
consolidation only \emph{moves} it between pools. For $J$ symmetric competitors the unit of
market rent splits three ways:
\[
  \underbrace{\tfrac{1}{J}}_{\text{protected baseline}}
  \;+\;
  \underbrace{\rho\,\tfrac{J-1}{J}}_{D:\ \text{dissipated by competition}}
  \;+\;
  \underbrace{(1-\rho)\,\tfrac{J-1}{J}}_{K:\ \text{curvature shield (variety)}}
  \;=\;1 .
\]
$D$ is rent competition wastes (redundant overhead, price wars); $K$ is surplus protected by
product variety/differentiation; $1/J$ is the monopoly baseline. The dial $\rho\in(0,1)$ is
the \emph{substitutes-vs-complements} regime: high $\rho$ = mostly wasteful competition,
low $\rho$ = mostly genuine variety. A roll-up runs two checks off this identity.

% ======================================================================
\section{Step 0 --- Define the market}
% ======================================================================
Fix the relevant antitrust/economic market and count the effective competitors $J$. Synergy
is measured \emph{within this market}; widen or narrow the boundary and the numbers move.
\emph{State it.}

% ======================================================================
\section{Check 1 --- Durable vs.\ fragile (the $\rho$ split)}
% ======================================================================

\subsection{The mechanics}
Consolidating one competitor, $J\to J-1$, moves rent into the protected baseline:
\[
  \Delta\!\left(\tfrac{1}{J}\right)=\tfrac{1}{J-1}-\tfrac1J=\frac{1}{J(J-1)}\quad(\text{captured}),
\]
sourced, because the total is conserved, as
\[
  \underbrace{\frac{\rho}{J(J-1)}}_{\text{recaptured waste }(-\Delta D)\ =\ \durable}
  \;+\;
  \underbrace{\frac{1-\rho}{J(J-1)}}_{\text{destroyed variety }(-\Delta K)\ =\ \fragile}.
\]
So of \emph{every} dollar of captured synergy, fraction $\rho$ is \durable{} (recaptured
waste --- a real efficiency gain that takes nothing from anyone) and fraction $(1-\rho)$ is
\fragile{} (variety converted to rent --- a transfer from customers, the antitrust-relevant
part, which invites entry and reverts). The split is $J$-independent; only its \emph{size}
$1/(J(J-1))$ depends on $J$.

\subsection{Estimating $\rho$ --- the one judgment that drives everything}
\begin{judgment}
$\rho$ is exogenous and it \emph{is} the analysis. Estimate it as the share of the projected
synergy that is genuine cost/efficiency recapture (consolidating duplicated functions,
ending value-destroying price competition) versus pure pricing power from fewer rivals. Anchor
it: scrutinize the synergy model line by line, classify each item as efficiency or
price-from-concentration, and let $\rho$ be the efficiency share. Document the call.
\end{judgment}

\subsection{Worked example}
A platform in a market of $J=6$ pitches consolidation to $3$, projecting \textbf{\$60M} of
EBITDA synergy --- which, being margin from consolidation, is captured rent in full. We
estimate the durable share at $\rho=0.60$ (60\% efficiency recapture, 40\% competition
reduction):
\begin{center}
\renewcommand{\arraystretch}{1.25}\small
\begin{tabular}{l r l}
\toprule
\textbf{Projected synergy} & \textbf{\$M} & \textbf{Character} \\
\midrule
Efficiency recaptured \;($\rho=0.60$)            & 36 & \durable{} --- sticks \\
Competition-reduction rent \;($1-\rho=0.40$)     & 24 & \fragile{} --- reverts \\
\midrule
\textbf{Captured synergy} & \textbf{60} & \\
\bottomrule
\end{tabular}
\end{center}
\textbf{\$36M durable; \$24M fragile.} Price the platform on the \$36M.

% ======================================================================
\section{Check 2 --- The entry floor (how far you can hold it)}
% ======================================================================

\subsection{The mechanics (Salop)}
A differentiated market with circular product space of differentiation $t$ and entry cost $f$
has equilibrium per-firm markup and free-entry count
\[
  \mu(n)=\frac{t}{n},\qquad
  n^{*}=\sqrt{t/f},\qquad
  \mu(n^{*})=\sqrt{tf}\ \ (\text{the entry-proof markup}),\qquad
  \text{durable HHI}=\frac{1}{n^{*}}=\sqrt{f/t}.
\]
Markup rises as players fall ($\mu=t/n$), and a merger lifts it by $\Delta\mu=t/(J(J-1))$. But
$n^{*}$ is the level entry restores: consolidating \emph{below} $n^{*}$ pushes markup above
$\sqrt{tf}$, and that excess is \fragile{} --- it draws entry back to $n^{*}$.

\subsection{Worked example}
With differentiation $t=\$20\text{M}$ and entry cost $f=\$1.25\text{M}$:
\[
  n^{*}=\sqrt{20/1.25}=\sqrt{16}=4\ \text{players},\quad
  \mu(n^{*})=\sqrt{20\cdot1.25}=\$5\text{M},\quad
  \text{durable HHI}=1/4=0.25 .
\]
\begin{center}
\renewcommand{\arraystretch}{1.25}\small
\begin{tabular}{l l}
\toprule
\textbf{Consolidation step} & \textbf{Character} \\
\midrule
$6\to4$ players & \durable{} --- down to the natural floor $n^{*}=4$ \\
$4\to3$ players & \fragile{} --- below $n^{*}$; entry restores $\sim$4 and competes the gain away \\
\bottomrule
\end{tabular}
\end{center}

\begin{note}
\textbf{The two checks are complementary, never added.} Check 1 sizes the fragile rent in
dollars (\$24M, via $\rho$); Check 2 locates it in the deal sequence (the $4\to3$ step, via
the floor). They are different lenses on the same fragile rent --- one in \$, one in deal
count --- and are \emph{not} summed. (Converting Check 2's per-unit markup into total dollars
would need a volume input, which is why they stay separate.)
\end{note}

% ======================================================================
\section{``How far should we consolidate?'' --- and why there is no forced optimum}
% ======================================================================
Netting the recaptured efficiency against the destroyed variety gives the per-merger value
(in the curvature metric)
\[
  \frac{\rho}{J(J-1)}-\frac{1-\rho}{J(J-1)}=\frac{2\rho-1}{J(J-1)},
  \qquad\text{cumulative }J\to n:\ (2\rho-1)\,\Delta\text{HHI}.
\]
The sign is fixed by $(2\rho-1)$ and the magnitude \emph{grows} as $J\to2$ --- increasing
returns, not diminishing. So the geometry alone gives a \textbf{corner solution}: consolidate
fully if $\rho>\tfrac12$, not at all if $\rho<\tfrac12$. \emph{There is no interior ``optimal
number of acquisitions'' in the math.} An interior stopping point appears only when you
\textbf{feed} a rising cost --- antitrust hazard or integration dis-synergy --- or impose the
entry floor of Check 2. The defensible recommendation is therefore the \emph{durable ceiling}
($n^{*}$), not a profit-maximizing $k$: in the example, $\rho=0.6>\tfrac12$ says ``keep going''
on pure rent, but the floor says everything below 4 players is fragile, so 4 is the durable
stopping point and the $4\to3$ deal buys rent that reverts.

% ======================================================================
\section{Judgment checkpoints (forced vs.\ fed)}
% ======================================================================
\begin{center}
\renewcommand{\arraystretch}{1.25}\small
\begin{tabular}{p{6.4cm} p{6.4cm}}
\toprule
\textbf{Forced by the algebra} & \textbf{Fed by you (state as assumptions)} \\
\midrule
The rent identity $1/J+D+K=1$; merger curvature loss $(1-\rho)/(J(J-1))$; HHI $=1/J+$ dispersion.
& $\rho$, the substitutes-vs-complements regime (drives the entire split). \\
Salop markup $t/n$; free-entry $n^{*}=\sqrt{t/f}$, markup $\sqrt{tf}$. & Differentiation $t$ and entry cost $f$ (locate the floor). \\
Capture-not-creation (the conserved total). & A price $w$ to put \$ on destroyed variety; the bridge that one market is both $\rho$- and $t$-shaped. \\
Net per-merger value $(2\rho-1)/(J(J-1))$. & Antitrust thresholds, integration cost (to get an interior $k$); the reversion \emph{timeline}. \\
\bottomrule
\end{tabular}
\end{center}
\textbf{Out of scope entirely:} the \emph{speed} of reversion. The test gives the durable
\emph{level} (a comparative-static fixed point), never how fast entry competes the fragile
rent away --- flag it as a judgment, never invent a number.

% ======================================================================
\section{Why it works (self-contained)}
% ======================================================================
Each step has an elementary reason --- a normalized partition, one line of algebra, and the
textbook free-entry condition. Nothing reaches beyond that.
\begin{itemize}
  \item \textbf{The rent decomposition is bookkeeping on a normalized total.} Set total market
  rent to $1$. With $J$ equal competitors the symmetric floor is $1/J$; the rest, $(J-1)/J$, is
  contested. Let $\rho$ be the fraction of that contested rent which is wasteful competition
  (recoverable) and $(1-\rho)$ the fraction protected by differentiation. The three pieces sum to
  $1$ by construction --- it's a partition, not an assumption. ($\rho$ is the one input you
  estimate.)
  \item \textbf{The merger move is one line.} $\Delta(1/J)=\tfrac{1}{J-1}-\tfrac1J=\tfrac{1}{J(J-1)}$
  shifts into the protected baseline, of which $\rho$ is durable (recovered waste) and $(1-\rho)$
  is fragile (lost variety). Net per merger $=\tfrac{2\rho-1}{J(J-1)}$ --- positive iff
  $\rho>\tfrac12$.
  \item \textbf{The entry floor is the free-entry condition.} In the standard circular-market
  model each firm's markup is $t/n$ and it serves $1/n$ of the market, so per-firm operating
  profit is $\tfrac{t}{n}\cdot\tfrac1n=\tfrac{t}{n^2}$. Free entry drives that to the entry cost
  $f$: $\tfrac{t}{n^2}=f\Rightarrow n^{*}=\sqrt{t/f}$, with entry-proof markup $t/n^{*}=\sqrt{tf}$
  and durable concentration $1/n^{*}=\sqrt{f/t}$. Four lines, standard IO.
\end{itemize}
\textbf{What's exact is the split and the floor formula; the inputs} ($\rho$, $t$, $f$)
\textbf{are judgment, and there is no ``optimal $k$''} --- the geometry gives a corner, so an
interior stop needs a fed cost (antitrust / integration). Present it as ``exact and
transparent,'' never as an external proof.

\vspace{6pt}\hrule\vspace{5pt}
{\footnotesize\itshape Companion to ``The Synergy Reality Test'' and to the Value-Creation
Bridge guide. Draft prepared by the ENT 105 team for Creo Advisors. Figures are illustrative
and reconcile with the one-pager. The rent split and the entry floor are grounded in
conservation-based accounting; $\rho$, $t$, $f$,
antitrust thresholds, and the reversion timeline are client-specific inputs set at scoping.}

\end{document}
